% Status: ueberarbeitet nach FHNW-Struktur fuer empirische Arbeiten
% (docs/project/FHNW_ACADEMICGUIDE_RULES.md, Abschnitt 4a + 10).
% Evidence: method_h1_brier_dm; method_h2_event_window; method_h3_wallet_tiers.
% Finale Zitation nach Quellenreview (docs/project/SOURCE_REVIEW_LOG.md).
\chapter{Einleitung}
\label{ch:intro}

\section{Ausgangslage und Problemstellung}
Prognosemaerkte buendeln das verteilte Wissen vieler Marktteilnehmender in einem
einzigen, laufend aktualisierten Preis, der sich als Wahrscheinlichkeit fuer ein
kuenftiges Ereignis lesen laesst. Die theoretische Grundlage bildet die
Markteffizienzhypothese, nach der Preise die verfuegbare Information widerspiegeln
\citep{lit_emh_001}. Im US-Wahljahr 2024 wurde der dezentrale, auf der Blockchain
abgewickelte Markt Polymarket zum sichtbarsten Beispiel dieser Idee. Er erreichte
zeitweise sehr hohe Handelsvolumina und stand wiederholt in Spannung zu
traditionellen Umfragen und Forecast-Modellen, die ein knapperes Rennen zeichneten
\citep{zotero_poly_001,zotero_poly_009}.

Diese Diskrepanz wirft eine grundlegende Frage auf. Verarbeitet ein solcher Markt
Information tatsaechlich besser als etablierte Quellen, oder erscheint er nur
deshalb ueberlegen, weil einzelne Episoden im Nachhinein betont werden? Die
vorliegende Arbeit nimmt diese Frage zum Ausgangspunkt und untersucht die
informationelle Effizienz von Polymarket im Kontext der US-Praesidentschaftswahl
2024 --- nicht als pauschale Behauptung, sondern als nachvollziehbar gepruefte,
scope-spezifische Aussage.

\section{Relevanz und Einordnung}
Die Fragestellung ist wirtschaftlich wie methodisch relevant. Erstens ruecken
Prognosemaerkte als kostenguenstige, laufend aktualisierte Forecast-Infrastruktur
in den Blick von Unternehmen, Medien und oeffentlichen Institutionen, die auf
verlaessliche Wahrscheinlichkeitsschaetzungen angewiesen sind. Zweitens macht die
Blockchain-Abwicklung jede Transaktion oeffentlich nachvollziehbar und eroeffnet
damit Analysen auf Wallet-Ebene, die in traditionellen Maerkten nicht moeglich
sind \citep{zotero_poly_001}. Drittens beruehrt die Diskussion um Grosskonten und
moeglichen informierten Handel Fragen der Marktintegritaet und der Regulierung
\citep{zotero_poly_005,zotero_poly_007}.

Eine reproduzierbare und methodisch vorsichtige Untersuchung leistet damit einen
doppelten Beitrag. Sie ordnet die empirische Effizienzforschung zu dezentralen
Maerkten ein und liefert zugleich eine praxisnahe Grundlage, um die Verlaesslichkeit
solcher Maerkte realistisch und ohne ueberzogene Erwartungen einzuschaetzen.

\section{Forschungsfrage und Hypothesen}
Aus der Problemstellung leitet sich die folgende Forschungsfrage ab:

\begin{quote}
In welchem Mass verarbeitet Polymarket waehrend der US-Praesidentschaftswahl 2024
Information, verglichen mit traditionellen Prognosequellen?
\end{quote}

Informationelle Effizienz \citep{lit_emh_001} wird dabei nicht direkt behauptet,
sondern ueber drei reproduzierbare, deterministisch pruefbare Proxy-Hypothesen
operationalisiert:

\begin{itemize}
  \item \textbf{H1 (Prognosequalitaet):} Polymarket liefert besser oder gleich gut
  kalibrierte Wahrscheinlichkeiten wie vergleichbare traditionelle Forecasts,
  gemessen am Brier-Verlust \citep{lit_brier_001} und am
  Diebold-Mariano-Vergleich der Verlustreihen \citep{lit_dm_001}.
  \item \textbf{H2 (Ereignisreaktion):} Die Polymarket-Wahrscheinlichkeiten bewegen
  sich um vorab kuratierte oeffentliche Ereignisse in plausibler Richtung,
  untersucht mit einem Ereignisstudien-Design \citep{lit_eventstudy_001}.
  \item \textbf{H3 (Wallet-Timing):} Dataset-relative Wallet-Tiers zeigen zeitliche
  Aktivitaetsmuster um Preisbewegungen, gemessen ueber Lead-Lag-Korrelationen und
  Granger-Tests \citep{lit_granger_001}.
\end{itemize}

Die drei Hypothesen greifen ineinander. H1 prueft die Ergebnisqualitaet der
Prognose, H2 die Reaktionsrichtung auf neue oeffentliche Information und H3 die
zeitliche Struktur der zugrunde liegenden Handelsaktivitaet. Zusammen erlauben sie
eine differenzierte statt einer pauschalen Antwort auf die Forschungsfrage.

\section{Abgrenzung des Untersuchungsgegenstands}
Der Untersuchungsgegenstand ist bewusst eng gefasst. Untersucht wird die US-Wahl
2024 als einzelner Fallkontext, Polymarket bildet die primaere Zeitreihe,
traditionelle Umfragen und Modell-Forecasts dienen als Vergleich. Die Arbeit erhebt
ausdruecklich keinen Anspruch auf universelle Markteffizienz, auf
Intraday-Reaktionsgeschwindigkeit, auf kausale Insider-Effekte oder auf handelbare
Profitabilitaet. Die Wallet-Analyse stuetzt sich auf oeffentliche, quellseitig
gefilterte Kauf-Transaktionen und ist deshalb als deskriptive Timing-Diagnostik zu
lesen, nicht als Nachweis privater Information.

Diese Selbstbeschraenkung folgt der in der Literatur betonten Vorsicht. Grosskonto-
Aktivitaet kann mit heterogenen Erwartungen statt mit Manipulation vereinbar sein
\citep{zotero_poly_001}, und Prognosemaerkte koennen reflexive Verzerrungen
aufweisen, die ihre Genauigkeit begrenzen \citep{zotero_poly_009}.

\section{Forschungsluecke und Beitrag}
Die bestehende Literatur liefert Theorie zur Informationsaggregation, etablierte
Messverfahren und gemischte Evidenz zu Polymarket, behandelt diese Aspekte aber
meist getrennt \citep{zotero_poly_002,zotero_poly_007}. Was fehlt, ist eine
reproduzierbare, deterministisch gepruefte Zusammenfuehrung von Prognosequalitaet,
Ereignisreaktion und On-Chain-Timing fuer denselben Marktkontext. Genau diese Luecke
schliesst die Arbeit fuer die US-Wahl 2024.

Der empirische Kern ist ein vollstaendig deterministischer, getesteter
Python-Analysepfad, der alle Kennzahlen berechnet, bevor sie interpretiert werden.
Als Sekundaerbeitraege dokumentiert die Arbeit ein aus der Untersuchung
hervorgegangenes read-only Monitoring-Werkzeug sowie die Multiagenten-Methodik
seiner Entwicklung (Kapitel~\ref{ch:erweiterungen}). Diese Sekundaerbeitraege
stuetzen die Forschungsfrage, ersetzen aber den empirischen Beleg nicht und
ueberlagern ihn nicht.

\section{Aufbau der Arbeit}
Die Arbeit ist in der Struktur einer empirischen Forschungsarbeit aufgebaut.
Kapitel~\ref{ch:theorie} ordnet Theorie und Literatur ein und entwickelt die
Forschungsluecke. Kapitel~\ref{ch:methodik} beschreibt Daten, Aufbereitung und die
deterministischen Verfahren fuer H1 bis H3. Die Kapitel~\ref{ch:h1} bis
\ref{ch:h3} berichten die Ergebnisse der drei Hypothesen. Kapitel~\ref{ch:erweiterungen}
stellt die Erweiterungen vor, darunter den Monitor, die Schweizer Fallstudie und
den praktischen Beitrag. Kapitel~\ref{ch:diskussion} fuehrt die Befunde zusammen,
diskutiert die Limitationen und zieht das Fazit, und Kapitel~\ref{ch:ausblick}
skizziert den agentengestuetzten Ausblick.
